% ============================================================================
% Chapter 8: Conclusion & Future Recommendations
% (formerly Chapter 10)
% ============================================================================

\chapter{Conclusion and Future Recommendations}
\label{chap:conclusion}

\section{Conclusion}
\label{sec:conclusion}

This graduation project successfully designed, developed, and implemented the \textbf{Autonomous Unified Swarm Robotics Architecture (AUSRA)} --- a comprehensive, cost-effective framework for multi-robot coordination and autonomous navigation. The project addressed the growing need for scalable swarm robotics platforms suitable for research and industrial applications including warehouse automation, disaster response, and search-and-rescue operations.

\subsection{Summary of Achievements}

\textbf{Mechanical Design:}
\begin{itemize}
    \item Developed a modular, three-layer chassis design optimized for omnidirectional mobility using three 58mm omni wheels in a 120° radial configuration
    \item Conducted systematic component selection using weighted decision matrices for motors, materials, and structural elements
    \item Validated structural integrity through Finite Element Analysis (FEA), achieving a safety factor exceeding 31 with 3mm aluminum plate configuration
    \item Achieved target payload capacity of 5.2 kg with extended battery runtime of 1.8 hours
    \item Created comprehensive URDF/Xacro robot description with Gazebo simulation integration
\end{itemize}

\textbf{Electrical System:}
\begin{itemize}
    \item Designed and implemented a robust power distribution network with dual-stage buck conversion (20V $\rightarrow$ 12V $\rightarrow$ 5V) providing isolated logic and actuator power rails
    \item Integrated comprehensive sensor suite: RPLIDAR A1 (360° LiDAR), OAK-D Lite depth camera with onboard AI, and MPU-9250 9-axis IMU
    \item Developed custom PCB shields for ESP32-S3 signal distribution and IMU interface with reliable screw terminal connections
    \item Implemented low-level firmware achieving 100 Hz PID motor control loops with encoder feedback at 600 PPR resolution
    \item Established micro-ROS communication over UART at 921600 baud with measured latency under 3 ms and packet loss below 0.02\%
\end{itemize}

\textbf{Autonomous Software Stack:}
\begin{itemize}
    \item Built complete navigation pipeline on ROS 2 Humble with Nav2 framework
    \item Implemented \texttt{slam\_toolbox} for asynchronous 2D SLAM achieving over 96\% mapping coverage across test environments
    \item Configured SmacPlanner2D (A*) for global path planning and DWB controller optimized for holonomic motion with full lateral velocity sampling
    \item Developed custom frontier-based exploration algorithm with BFS clustering and 98\% safety threshold filtering, achieving 97\% coverage of $20 \times 20$ m environments in under 3 minutes
    \item Implemented omnidirectional kinematic model enabling heading-independent translation with accurate odometry covariance propagation for EKF fusion
\end{itemize}

\textbf{System Integration:}
\begin{itemize}
    \item Successfully integrated all subsystems into functional prototypes with Docker containerization using NVIDIA runtime for reproducible deployment
    \item Demonstrated multi-robot support through namespace isolation in TF tree and topic structure (e.g., \texttt{ausra\_1}, \texttt{ausra\_2})
    \item Achieved navigation success rates between 92--100\% across scenarios including static obstacles, narrow passages (0.6m), and dynamic environments
    \item Validated complete autonomous stack achieving 95\% map coverage in 1 minute 40 seconds with 100\% goal success rate (17/17 navigation goals)
    \item Maintained real-time performance with total system resource utilization well within Jetson Orin Nano capabilities ($<$50\% CPU, $<$2.5 GB RAM)
\end{itemize}

\subsection{Contributions}
This project makes the following contributions:
\begin{enumerate}
    \item A unified, modular robot platform at \textbf{\$420 USD per unit}---23.6\% cost reduction compared to TurtleBot3 while offering superior capabilities including omnidirectional motion, GPU-accelerated AI processing, and native swarm support
    \item Comprehensive documentation serving as a reference for future swarm robotics research projects
    \item Open-source software stack built on ROS 2 Humble with Nav2, including custom frontier exploration and omnidirectional driver packages
    \item Validated design methodology using selection matrices and systematic analysis for component choices (IMU, SLAM algorithm, compute platform, depth camera)
    \item Complete simulation environment with Gazebo Classic worlds and sensor plugins for hardware-in-the-loop testing before physical deployment
\end{enumerate}

\subsection{Project Metrics Summary}

\begin{table}[H]
\centering
\caption{Final Project Achievement Summary}
\label{tab:conclusion_metrics}
\begin{tabular}{|l|c|c|}
\hline
\textbf{Metric} & \textbf{Target} & \textbf{Achieved} \\
\hline
Cost per Robot & $<$ \$500 & \$420 \\
SLAM Coverage & $>$ 90\% & 96\%+ \\
Navigation Success & $>$ 90\% & 92--95\% \\
Exploration Coverage & $>$ 90\% & 95\% \\
Battery Runtime & $>$ 1.0 hrs & 1.25 hrs \\
Maximum Speed & 0.5 m/s & 0.52 m/s \\

\hline
\end{tabular}
\end{table}

\section{Future Recommendations}
\label{sec:future_work}

While the AUSRA project achieved its primary objectives, several areas present opportunities for future development:

\subsection{Hardware Improvements}
\begin{itemize}
    \item \textbf{Direct IMU Integration:} Complete transition from ESP32-bridged IMU to direct Jetson I2C connection
    \item \textbf{Docking Charging:} Implementation of autonomous docking stations with station charging for extended autonomous operation
    \item \textbf{Improved Encoders:} Upgrade to higher resolution encoders ($>$1000 PPR) for improved odometry accuracy in precision applications
\end{itemize}

\subsection{Software Enhancements}
\begin{itemize}
    \item \textbf{Single Pose Graph SLAM:} Implement unified multi-robot SLAM where all agents contribute to a shared constraint graph, enabling loop closures from one robot to correct drift in others
    \item \textbf{mTSP Task Allocation:} Integrate Multiple Traveling Salesman Problem solver (using Google OR-Tools) for globally optimal frontier assignment, replacing greedy nearest-frontier selection
    \item \textbf{Velocity Obstacle Avoidance:} Implement decentralized collision avoidance using velocity obstacles where agents broadcast velocity vectors to neighbors for real-time trajectory deconfliction
    \item \textbf{OAK-D Vision Pipeline:} Utilize the depth camera's onboard AI capabilities for AprilTag detection (victim identification) and semantic segmentation via DepthAI SDK
\end{itemize}

\subsection{Scalability Testing}
\begin{itemize}
    \item \textbf{Larger Swarms(in simualtion0):} Systematic testing with 5--10+ robots to validate namespace isolation scalability and network bandwidth requirements
    \item \textbf{Distributed Computing:} Offload computationally intensive tasks (e.g., global map merging) to edge servers while maintaining local autonomy
\end{itemize}

\subsection{Real-World Deployment Considerations}
\begin{itemize}
    \item \textbf{Outdoor Operation:} Ruggedization for outdoor localization
    \item \textbf{Dynamic Environments:} Enhanced handling of moving obstacles (humans, forklifts) through predictive trajectory modeling
    \item \textbf{Long-Duration Missions:} Battery hot-swap mechanisms and mission resumption after charging cycles
\end{itemize}

\subsection{Research Directions}
\begin{itemize}
    \item \textbf{Reinforcement Learning:} Train navigation policies using simulation-to-real transfer for improved performance in cluttered environments
    \item \textbf{Swarm Intelligence:} Investigate bio-inspired coordination algorithms (ant colony optimization, particle swarm) for emergent collective behavior
    \item \textbf{Human-Swarm Interaction:} Develop intuitive interfaces allowing single operators to supervise and direct multiple robots through high-level commands
\end{itemize}
\newpage
\vspace{1cm}
\hrule
\vspace{0.5cm}

The AUSRA project demonstrates the feasibility of building a cost-effective, unified swarm robotics platform using commercially available components and open-source software. By achieving all primary design objectives---autonomous navigation, 2D SLAM mapping, multi-robot coordination, and sub-\$500 unit cost---AUSRA provides a validated foundation for collaborative robotics research. The modular architecture, comprehensive documentation, and ROS 2 integration enable continued development toward practical applications in warehouse automation, agricultural monitoring, and search-and-rescue operations. Future work focusing on advanced multi-robot coordination algorithms, enhanced perception capabilities, and real-world deployment considerations will further extend the platform's capabilities and impact.
